\subsection{Riemann Integration}

The integration side of one-variable calculus is organized by Riemann and Riemann--Stieltjes sums; the fundamental theorem connects the two halves.

This subsection develops Riemann and Riemann--Stieltjes integration by partitions, then links integrals to antiderivatives through the Fundamental Theorem of Calculus. Point partitions and interval-envelope partitions are equivalent via common refinement.

\begin{definition} \label{def:partition-and-riemann-upper-lower-sums}\label{def:riemann-stieltjes-core-source-todo} \hlt{Partitions and Riemann Upper/Lower Sums}\\
A \hlt{partition} of $[a,b]$ is $\Pi=\{x_0,\ldots,x_n\}$ with $a=x_0\leq x_1\leq\cdots\leq x_n=b$, $\Delta x_i\defeq x_i-x_{i-1}$.\\
For bounded real $f$ on $[a,b]$, set $M_i\defeq\sup_{x_{i-1}\leq x\leq x_i}f(x)$, $m_i\defeq\inf_{x_{i-1}\leq x\leq x_i}f(x)$,\\
and $U(\Pi,f)\defeq\sum\limits_{i=1}^{n}M_i\Delta x_i$, $L(\Pi,f)\defeq\sum\limits_{i=1}^{n}m_i\Delta x_i$.\\
Define $\overline{\int_a^b f(x)\,dx}\defeq\inf_\Pi U(\Pi,f)$, $\underline{\int_a^b f(x)\,dx}\defeq\sup_\Pi L(\Pi,f)$, with $\Pi$ ranging over all partitions.\\
If these agree, $f$ is \hlt{Riemann integrable}, written $f\in\mathcal{R}$, and the common value is $\int_a^b f(x)\,dx$ or $\int_a^b f$.\\
For $m\leq f\leq M$, $m(b-a)\leq L(\Pi,f)\leq U(\Pi,f)\leq M(b-a)$, upper and lower integrals exist for bounded $f$.
\end{definition}

\begin{figure}[H]
  \centering
  \begin{tikzpicture}[
  x=0.80cm,
  y=0.52cm,
  >=Stealth,
  every node/.style={font=\scriptsize},
  axis/.style={->, gray!70, line width=0.45pt},
  curve/.style={green!45!black, line width=0.8pt},
  rect/.style={draw=black!70, fill=gray!5, line width=0.55pt},
  tag/.style={circle, fill=black, inner sep=1.25pt},
  guide/.style={gray!55, dashed, dash pattern=on 2.5pt off 2pt, line width=0.45pt}
]
  \draw[axis] (-0.45,0) -- (7.75,0) node[right] {$x$};
  \draw[axis] (0,-0.55) -- (0,2.85) node[above] {$y$};

  \coordinate (COneTop) at (1.45,1.08);
  \coordinate (CTwoTop) at (2.50,2.48);
  \coordinate (CThreeTop) at (3.55,1.52);
  \coordinate (CnTop) at (6.70,0.78);

  \foreach \xl/\xr/\c/\h in {
    1.05/2.05/1.45/1.08,
    2.05/3.05/2.50/2.48,
    3.05/4.05/3.55/1.52,
    6.20/7.20/6.70/0.78
  } {
    \draw[rect] (\xl,0) rectangle (\xr,\h);
    \draw[guide] (\c,0) -- (\c,\h);
    \node[tag] at (\c,0) {};
  }

  \draw[curve]
    (1.05,0.74) .. controls (1.16,0.74) and (1.30,0.78) .. (COneTop)
    .. controls (1.62,1.42) and (2.06,2.50) .. (CTwoTop)
    .. controls (2.90,2.38) and (3.20,1.82) .. (CThreeTop)
    .. controls (4.38,0.62) and (5.05,0.42) .. (5.78,0.68)
    .. controls (6.15,0.82) and (6.40,0.84) .. (CnTop)
    .. controls (6.98,0.75) and (7.18,0.68) .. (7.38,0.72);

  \foreach \c/\lab in {1.45/c_1,2.50/c_2,3.55/c_3,6.70/c_n} {
    \node[tag] at (\c,0) {};
    \node[above=2pt] at (\c,0) {$\lab$};
  }

  \node[below=12pt] at (1.05,0) {$x_0=a$};
  \node[below=12pt] at (2.05,0) {$x_1$};
  \node[below=12pt] at (3.05,0) {$x_2$};
  \node[below=12pt] at (4.05,0) {$x_3$};
  \node[below=12pt] at (5.10,0) {$\cdots$};
  \node[below=12pt] at (6.20,0) {$x_{n-1}$};
  \node[below=12pt] at (7.20,0) {$x_n=b$};
\end{tikzpicture}

  \caption{A tagged Riemann sum uses one sample value $f(c_i)$ on each subinterval $[x_{i-1},x_i]$.}
  \label{fig:riemann-tagged-sum}
\end{figure}

\begin{definition} \label{def:tagged-partitions-and-riemann-sums} \hlt{Tagged Partitions and Riemann Sums}\\
A \hlt{tagged partition} of $[a,b]$ is a partition $\Pi=\{x_0,\ldots,x_n\}$ together with tags $c_i\in[x_{i-1},x_i]$.\\
For the ordinary Riemann integral, the corresponding \hlt{Riemann sum} is $R(f,\Pi,\{c_i\})\defeq\sum\limits_{i=1}^{n}f(c_i)\Delta x_i$.\\
For every tagged partition, $L(\Pi,f)\leq R(f,\Pi,\{c_i\})\leq U(\Pi,f)$.
\end{definition}

\begin{perspective} \label{per:gauge-integral-context} \hlt{Tagged Sums and Gauges}\\
Gauge integration keeps tagged sums but replaces one mesh bound $\delta$ with a local bound $\delta(c_i)$.\\
Thus fineness may vary by tag, giving a stronger framework for FTC-type results.
\end{perspective}

\begin{figure}[H]
  \centering
  \begin{tikzpicture}[
  scale=0.72,
  >=Stealth,
  every node/.style={font=\scriptsize},
  axis/.style={->, gray!70, line width=0.45pt},
  curve/.style={blue!70!black, line width=0.9pt},
  upper/.style={fill=orange!18, draw=orange!75!black, line width=0.45pt},
  lower/.style={fill=teal!18, draw=teal!70!black, line width=0.45pt},
  guide/.style={gray!45, dashed, line width=0.35pt}
]
  \draw[axis] (-0.35,0) -- (7.8,0) node[right] {$x$};
  \draw[axis] (0,-0.2) -- (0,3.9) node[above] {$y$};

  \foreach \x/\lab in {0/$a$,1.45/$x_1$,3.2/$x_2$,5.15/$x_3$,7.0/$b$} {
    \draw[guide] (\x,0) -- (\x,3.55);
    \node[below] at (\x,0) {\lab};
  }

  \foreach \xl/\xr/\M/\m in {
    0/1.45/2.58/1.82,
    1.45/3.2/1.82/0.83,
    3.2/5.15/2.25/0.80,
    5.15/7.0/1.75/0.61
  } {
    \draw[upper] (\xl,0) rectangle (\xr,\M);
    \draw[lower] (\xl,0) rectangle (\xr,\m);
  }

  \draw[curve, domain=0:7, samples=120]
    plot (\x,{1.85 + 0.55*sin(90*\x) + 0.35*sin(185*\x) - 0.07*\x});

  \node[orange!75!black, anchor=west] at (7.15,2.55) {$U(P,f)$};
  \node[teal!70!black, anchor=west] at (7.15,1.0) {$L(P,f)$};
\end{tikzpicture}

  \caption{Upper and lower sums record the oscillation of $f$ on each partition interval.}
  \label{fig:riemann-upper-lower-partition}
\end{figure}

